
\section*{1. Formal Varieties over a Pseudocompact Ring}
\phantomsection
\label{sga3:VIIB:1}
\addcontentsline{toc}{section}{1. Formal Varieties over a Pseudocompact Ring}

\medskip
\noindent\textbf{1.1.}
\phantomsection
\label{sga3:VIIB:1.1}
To every pseudocompact ring \(A\) one may associate a formal scheme
(EGA~I, 10.4.2) as follows. Its underlying topological space is the
set \(\Upsilon(A)\) of open prime ideals of \(A\), which are therefore
maximal, endowed with the discrete topology. For every subset
\(E\subset\Upsilon(A)\), the space of sections of the structure sheaf
over \(E\) is the Cartesian product
\[
  \prod_{\mathfrak m\in E}A_{\mathfrak m}.
\]
The formal scheme so obtained is denoted by \(\operatorname{Spf}(A)\),
the \emph{formal spectrum} of \(A\).

If \(A\) and \(B\) are two pseudocompact rings, a \emph{morphism} from
\(\operatorname{Spf}(B)\) to \(\operatorname{Spf}(A)\) consists of a
map
\[
  f:\Upsilon(B)\longrightarrow\Upsilon(A)
\]
and a family of continuous homomorphisms
\[
  f_y^\natural:A_{f(y)}\longrightarrow B_y,
  \qquad y\in\Upsilon(B).
\]
Such a morphism induces a continuous homomorphism
\[
  f^\natural:
  A=\prod_{x\in\Upsilon(A)}A_x
  \longrightarrow
  B=\prod_{y\in\Upsilon(B)}B_y.
\]
The converse is true:

\medskip
\noindent\textbf{Proposition.}
\phantomsection
\label{sga3:VIIB:1.1-spf-fully-faithful}
\textit{The contravariant functor
\(A\mapsto\operatorname{Spf}(A)\) is fully faithful.}

Indeed, let \(\phi:A\to B\) be a continuous homomorphism. The inverse
image \(\phi^{-1}(\mathfrak n)\) of an open maximal ideal of \(B\) is
an open prime ideal of \(A\), and hence is maximal in \(A\). We
therefore obtain a map
\[
  \mathfrak n\longmapsto\phi^{-1}(\mathfrak n)
\]
from \(\Upsilon(B)\) to \(\Upsilon(A)\), and \(\phi\) induces a
continuous homomorphism
\[
  A_{\phi^{-1}(\mathfrak n)}\longrightarrow B_{\mathfrak n}.
\]
Thus \(\phi\) induces a morphism of formal schemes
\[
  \operatorname{Spf}(\phi):
  \operatorname{Spf}(B)\longrightarrow\operatorname{Spf}(A).
\]
It is readily verified that
\[
  \operatorname{Spf}(\phi)^\natural=\phi,
  \qquad
  \operatorname{Spf}(f^\natural)=f
\]
for every morphism
\(f:\operatorname{Spf}(B)\to\operatorname{Spf}(A)\), which proves the
proposition.

Although we shall speak here only of formal schemes of the form
\(\operatorname{Spf}(A)\), we shall use the language of formal schemes
rather than that of pseudocompact rings, so as to ground our assertions
in geometric intuition.

\medskip
\noindent\textbf{1.2.}
\phantomsection
\label{sga3:VIIB:1.2}
Let \(k\) be a pseudocompact ring.

\medskip
\noindent\textbf{Definition 1.2.A.}
\phantomsection
\label{sga3:VIIB:1.2.A}
A \emph{formal variety over \(k\)} is a formal scheme \(X\)
over \(\operatorname{Spf}(k)\) which is isomorphic to a formal
\(k\)-scheme \(\operatorname{Spf}(A)\) for some profinite
\(k\)-algebra \(A\). The algebra \(A\) is then isomorphic to the
\emph{affine algebra} of \(X\), that is, the algebra of sections of
the structure sheaf \(\mathcal O_X\) of \(X\).

We denote by \(\operatorname{Vaf}/k\) the full subcategory of the
category of formal schemes over \(\operatorname{Spf}(k)\) whose
objects are the formal \(k\)-varieties.

By \SGARefLink{sga3:VIIB:1.1}{1.1}, the functor \(A\mapsto\operatorname{Spf}(A)\) is an
anti-equivalence from \(\operatorname{Alp}/k\) (\SGARefLink{sga3:VIIB:0.4.1}{0.4.1}) onto
\(\operatorname{Vaf}/k\). Hence, by the corollary of \SGARefLink{sga3:VIIB:0.4.2}{0.4.2}:

\medskip
\noindent\textbf{Proposition 1.2.B.}
\phantomsection
\label{sga3:VIIB:1.2.B}
\textit{The category \(\operatorname{Vaf}/k\) has projective and
direct limits.}

For example, let \(X\in\operatorname{Ob}(\operatorname{Vaf}/k)\), let
\[
  f:Y\longrightarrow X,
  \qquad
  g:Z\longrightarrow X
\]
be two formal \(k\)-varieties over \(X\), and let \(A,B,C\) be the
affine algebras of \(X,Y,Z\), respectively.
By \SGARefLink{sga3:VIIB:0.4.1}{0.4.1}, the affine algebra of the fiber product \(Y\times_X Z\)
identifies with \(B\widehat{\otimes}_A C\),
Consequently, the inclusion of \(\operatorname{Vaf}/k\) into the
category of all formal \(k\)-schemes commutes with finite projective
limits (cf.\ EGA~I, 10.7).

Direct limits of formal \(k\)-varieties correspond to projective
limits of their affine algebras.

\medskip
\noindent\textbf{Example 1.2.C (Cokernels).}
\phantomsection
\label{sga3:VIIB:1.2.C}
Let \(d,e:X\rightrightarrows Y\) be a double arrow in
\(\operatorname{Vaf}/k\). The affine algebra of
\(\operatorname{Coker}(d,e)\) is isomorphic to the kernel of the
homomorphisms induced on the affine algebras of \(X\) and \(Y\).
There is also the following construction of
\(\operatorname{Coker}(d,e)\). Its underlying topological space is the
cokernel of the underlying spaces. Let \(p\) be the
canonical projection from the set underlying \(Y\) to this cokernel,
and let \(z\) be a point of the cokernel. The local algebra of
\(\operatorname{Coker}(d,e)\) at \(z\) is the kernel of the double
arrow
\begin{center}
  \phantomsection
  \label{sga3:VIIB:1.2.C-cokernel-double-arrow}
  \(
d^\natural,e^\natural:\quad
\prod_{p(y)=z}\mathcal O_{Y,y}
\mathrel{\substack{\longrightarrow\\[-.65ex]\longrightarrow}}
\prod_{q(x)=z}\mathcal O_{X,x}
\)

\end{center}
where
\[
  q=p\circ d=p\circ e,
\]
and \(d^\natural\) and \(e^\natural\) are induced by the
homomorphisms \(d_x^\natural\) and \(e_x^\natural\), in the notation
of \SGARefLink{sga3:VIIB:1.1}{1.1}.

\medskip
\noindent\textbf{Definition 1.2.D.}
\phantomsection
\label{sga3:VIIB:1.2.D}
Let \(\phi:k\to\ell\) be a homomorphism of pseudocompact
rings and let \(X\) be a formal \(k\)-variety with affine algebra
\(A\). The formal scheme
\[
  X\times_{\operatorname{Spf}(k)}\operatorname{Spf}(\ell),
\]
obtained by base change, is a formal \(\ell\)-variety. We also denote
it by
\[
  X\widehat{\otimes}_k\ell.
\]
Its affine algebra is the completed tensor product
\(A\widehat{\otimes}_k\ell\) (cf.\ 0.5 and EGA~I, \S~10).

\medskip
\noindent\textbf{Remark 1.2.E.}
\phantomsection
\label{sga3:VIIB:1.2.E}
Since every formal variety over \(k\) decomposes into formal
varieties over the local components of \(k\), \emph{in certain proofs
we shall assume that \(k\) is a local pseudocompact ring}.

We now give some examples while fixing our terminology.

\medskip
\noindent\textbf{1.2.1.}
\phantomsection
\label{sga3:VIIB:1.2.1}
A \(k\)-\emph{functor} is, by definition, a covariant functor from
\(\operatorname{Alf}/k\) to \((\mathbf{Sets})\). By \SGARefLink{sga3:VIIB:1.1}{1.1} and \SGARefLink{sga3:VIIB:0.4.2}{0.4.2},
we may identify
\[
  \operatorname{Vaf}/k
  \simeq
  (\operatorname{Alp}/k)^\circ
\]
with a full subcategory of the category of \(k\)-functors by
associating to every formal \(k\)-variety \(X\) the functor
\[
  \operatorname{Alf}/k\longrightarrow(\mathbf{Sets}),
  \qquad
  C\longmapsto
  X(C)=
  \operatorname{Hom}_{\operatorname{Vaf}/k}
    \bigl(\operatorname{Spf}(C),X\bigr).
\]

We shall later encounter \(k\)-functors \(\mathbf F\) which associate
to every object \(C\) of \(\operatorname{Alf}/k\) a \(C\)-module
\(\mathbf F(C)\), and to every morphism
\(\phi:C\to D\) in \(\operatorname{Alf}/k\) a \(k\)-linear map
\[
  \mathbf F(\phi):\mathbf F(C)\longrightarrow\mathbf F(D)
\]
such that, for \(x\in\mathbf F(C)\) and \(\lambda\in C\),
\[
  \mathbf F(\phi)(\lambda x)
  =
  \phi(\lambda)\mathbf F(\phi)(x).
\]
By \SGARefLink{sga3:i:definition:3-1}{Expos\'e I, 3.1}, such an \(\mathbf F\) is endowed with the structure
of an \(\mathbf O_k\)-module, where \(\mathbf O_k\) denotes the
\(k\)-functor in rings which assigns to every object \(C\) of
\(\operatorname{Alf}/k\) its underlying ring.

\medskip
\noindent\textbf{Definitions.}
\phantomsection
\label{sga3:VIIB:1.2.1-definitions}

\noindent\textit{\SGARefLink{sga3:vii:paragraph:5-1:equation:i}{(i)}}
An \(\mathbf O_k\)-module \(\mathbf F\) is called
\emph{admissible} if every morphism \(\phi:C\to D\) in
\(\operatorname{Alf}/k\) induces a bijection from
\(D\otimes_C\mathbf F(C)\) onto \(\mathbf F(D)\).

\smallskip
\noindent\textit{\SGARefLink{sga3:vii:paragraph:5-1:equation:ii}{(ii)}}
The module \(\mathbf F\) is called \emph{flat} if it is admissible and
if, for every object \(C\) of \(\operatorname{Alf}/k\), the
\(C\)-module \(\mathbf F(C)\) is flat.

For example, if \(M\) is a \(k\)-module, not necessarily
pseudocompact, we denote by \(\mathbf W(M)\), as in \SGARefLink{sga3:i:definition:4-6-1}{Expos\'e I, 4.6.1},
the \(\mathbf O_k\)-module
\[
  C\longmapsto C\otimes_kM.
\]
Then \(\mathbf W(M)\) is flat whenever \(M\) is flat over \(k\).

Moreover, the functor \(M\mapsto\mathbf W(M)\) has as right adjoint
the functor which assigns to every \(\mathbf O_k\)-module
\(\mathbf F\) the \(k\)-module
\[
  \varprojlim_{\mathfrak l}\mathbf F(k/\mathfrak l),
\]
where \(\mathfrak l\) ranges over the open ideals of \(k\).

\medskip
\noindent\textbf{1.2.2.}
\phantomsection
\label{sga3:VIIB:1.2.2}
In what follows, an \(\mathbf O_k\)-module will always be denoted by a
boldface letter such as \(\mathbf F\). When \(k\) is Artinian, we shall
simply write \(F\) instead of \(\mathbf F(k)\). In that case, it is
clear that the functor
\[
  \mathbf F\longmapsto F
\]
induces an equivalence from the category of flat
\(\mathbf O_k\)-modules onto the category of flat \(k\)-modules! The
purpose of this terminology is solely to allow us to reason ``as if
\(k\) were always Artinian.''

In accordance with Expos\'e I, \S~\SGARefLink{sga3:i:definition:3-1}{3.1}, we shall use analogous
conventions for other algebraic structures. Thus an
\(\mathbf O_k\)-algebra, respectively an \(\mathbf O_k\)-coalgebra,
an \(\mathbf O_k\)-Lie algebra, or an \(\mathbf O_k\)-\(p\)-Lie
algebra, consists of an \(\mathbf O_k\)-module \(\mathbf M\) together
with, for every object \(C\) of \(\operatorname{Alf}/k\), a structure
of \(C\)-algebra, respectively \(C\)-coalgebra, \(C\)-Lie algebra, or
\(C\)-\(p\)-Lie algebra, on \(\mathbf M(C)\). We further require that
for every morphism \(\phi:C\to D\) in \(\operatorname{Alf}/k\), the
map
\[
  D\otimes_C\mathbf M(C)\longrightarrow\mathbf M(D)
\]
induced by \(\mathbf M(\phi)\) be a homomorphism of \(D\)-algebras,
respectively of \(D\)-coalgebras, and so on.

Finally, if \(\mathbf F\) and \(\mathbf G\) are two
\(\mathbf O_k\)-modules, then \(\mathbf F\otimes_k\mathbf G\) denotes
the \(\mathbf O_k\)-module
\[
  C\longmapsto\mathbf F(C)\otimes_C\mathbf G(C).
\]

